\documentclass[11pt,a4paper]{article}
\usepackage[margin=2.4cm]{geometry}
\usepackage{parskip}
\usepackage[hidelinks]{hyperref}
\usepackage{xcolor}
\usepackage{booktabs}

% Anonymised: no author, affiliation, or location appears anywhere in this file.
\pagestyle{plain}
\newcommand{\point}[1]{\medskip\noindent\textbf{\color{black!75}#1}\par\smallskip}
\newcommand{\reply}{\noindent\textbf{Response.}\ }

\title{\vspace{-1.5cm}Response to the Editor\\[2pt]
\large Carbon Footprint and Carbon-Aware Selection of Document-Oriented,\\
Relational, and Hybrid Object-Relational Storage for Image-Based\\
Time-Series Workloads in Green IoT}
\date{}

\begin{document}
\maketitle
\vspace{-1.2cm}

\noindent We thank the Editor for a constructive and precise assessment, and for the opportunity to
revise. Every point has been addressed, and the manuscript is substantially stronger as a result. In
particular, the request for stronger evidence of applicability prompted a new set of experiments
against real recorded camera payloads, which both confirmed the paper's central conclusion and
uncovered a quantitative correction we would otherwise have reported uncorrected. We are grateful
for that.

Below, each of the Editor's points is quoted, followed by our response and the location of the
change. One point is answered with partial rebuttal, which we set out explicitly rather than leave
implicit.

\section*{1. Positioning within the existing literature}

\point{``The discussion should better position the proposed framework within the existing scientific
literature, demonstrating more convincingly how it advances beyond previous studies rather than
primarily summarizing the experimental findings.''}

\reply We accept this fully. On re-reading, the Discussion explained \emph{mechanism} and then moved
directly to threats to validity, without returning to the literature at all. Every citation in that
section was to vendor or technical documentation rather than to scholarly work. The Editor's
characterisation was accurate.

We have added a new Discussion subsection, \textbf{``Positioning against prior work''}, which states
what changes relative to prior studies rather than restating what we measured:

\begin{itemize}
  \item \textbf{An accepted engine ordering is inverted.} Prior benchmarks on scalar time-series data
  report MongoDB ingesting faster than TimescaleDB. We reproduce that at low resolution
  ($>6\times$ faster at 360p) and then show it does not survive its own premise: the ordering
  inverts between 1440p and 4K, and MongoDB becomes \emph{infeasible} at 6K. What the scalar
  literature presents as a property of the engines is a property of the payload regime those
  benchmarks occupy.
  \item \textbf{The carbon-accounting literature supplies methods, not answers at this granularity.}
  Green Algorithms, the Software Carbon Intensity specification, and the embodied-carbon literature
  can attribute a footprint to a workload once its energy is known, but cannot indicate which
  architecture will consume that energy. We measure per architecture, per operation, and per
  resolution, and we note that a single aggregate figure per architecture, which those methods
  would ordinarily yield, would have averaged the crossover away and produced a spurious universal
  winner.
  \item \textbf{The output is a decision instrument, not a benchmark table.} The framework is now
  expressed as a transferable rule (see point 3) rather than a recommendation tied to our hardware.
\end{itemize}

We also state explicitly what we do \emph{not} claim: that the specific crossover resolutions
generalise. Table~1 has additionally been extended with a column recording what each prior study
concludes and what changes in our payload regime, so the advance is visible in the comparison itself
rather than asserted in prose.

\section*{2. Dialogue with recent publications in this journal}

\point{``The manuscript would also benefit from a stronger dialogue with recent publications in
Environment Systems and Decisions, reinforcing its contribution to the journal's body of
knowledge.''}

\reply Accepted. The submitted version cited no work from this journal. We have added a new Related
Work subsection engaging six recent articles substantively. Each is used to sharpen a specific
methodological choice, not cited in passing:

\begin{itemize}
  \item A systematic review of multi-criteria decision making for material selection identifies a
  recurring weakness: decision criteria are often taken from the literature rather than measured in
  the adopting organisation's own context. This directly motivates our decision to make the decision
  variable (payload size) measurable by the adopter, and to express the rule as a computed crossover
  rather than a fixed recommendation.
  \item A life-cycle assessment of lightweight railway vehicles establishes sensitivity analysis as
  what separates a robust conclusion from an artefact of assumed parameters. We identify our
  multi-region grid-intensity analysis as playing exactly that role, and show our ranking is
  invariant to the most uncertain parameter in any carbon accounting.
  \item A data-driven study of smart-home energy consumption treats building instrumentation as the
  object of environmental analysis. We draw the distinction that its models account for the energy
  the building consumes, but not for the tier that stores what the instrumentation records, which
  we show can be halved by a storage decision alone.
  \item Two studies of IoT operational integrity (self-healing via log analytics; real-time risk at
  the edge) both presuppose a data tier that durably retains device output. We supply the missing
  environmental cost of that presupposition.
  \item A review of price-driven energy management and demand response prompts a substantive point:
  because emissions are energy multiplied by grid intensity, temporal or geographic load shifting
  rescales every architecture equally and cannot change which one to select. The saving must come
  from the architecture itself.
\end{itemize}

\section*{3. Practical and scientific value, and applicability beyond our environment}

\point{``At present, the framework is derived almost entirely from the experimental environment
developed by the authors, and its applicability beyond the evaluated hardware, software
configuration, and workload characteristics remains insufficiently discussed.''}

\reply This was the most substantive point, and we have addressed it in three ways: analytically,
empirically with new experiments, and by separating claims that depend on measurement from those
that do not. We also offer one partial rebuttal, stated at the end.

\subsection*{3a. An analytical model, so the rule transfers rather than being copied}

We now derive the mechanism behind the crossover explicitly. Each architecture pays a \emph{fixed}
per-operation cost $a$ and a \emph{per-byte} cost $b$, so storing a frame of size $s$ costs
approximately $E(s) = a + bs$. Two architectures with different slopes cross exactly once, at
$s^{*} = (a_2-a_1)/(b_1-b_2)$.

MongoDB has the smallest fixed cost and the largest per-byte cost; the hybrid is its mirror image;
inline PostgreSQL lies between on both axes, which is why it is best at no operation and worst at
none. We fit these constants to our measurements and report them. \textbf{The consequence is that an
operator need not adopt our resolutions:} two insertion runs per architecture at two payload sizes
yield $a$ and $b$ on their hardware, and the equation gives their crossover. Our Table~5 is thus one
instantiation of the equation, and the equation is what generalises.

\subsection*{3b. New experiments: real camera payloads, a second camera, and a different host}

We describe these in the manuscript as \textbf{additional experiments carried out for this revision}.

The submitted benchmark inserted one synthetic collage for every row, so all 2{,}000 stored rows were
byte-identical, which is reproducible but not what a camera produces. We recorded a corpus of real frames
from an operational deployment (one networked camera, two USB cameras; ten minutes at 5 frames per
second; 2{,}845--2{,}988 frames per camera), verified every frame byte-distinct, and re-ran the
benchmark with a different real image for every inserted row. Three results follow:

\begin{enumerate}
  \item \textbf{The architectural conclusion holds.} MongoDB is the lowest-carbon architecture at
  every resolution from 360p to 1080p, as the synthetic sweep reports for that range, and the
  result replicates independently on a second camera.
  \item \textbf{The synthetic payload is close to a worst case,} not a typical one: real frames are
  $3.1\times$ to $5.4\times$ smaller at equal resolution.
  \item \textbf{The identical-row construction overstated MongoDB's margin.} MongoDB's storage
  amplification at 1080p rises from $0.10\times$ with identical rows to $0.92\times$ with distinct
  frames, because bucketed duplicates are removed almost entirely by compression while real frames
  are not. At matched payload size, inline PostgreSQL's carbon penalty relative to MongoDB falls
  from $3.96\times$ to $1.11\times$.
\end{enumerate}

We report this third finding prominently and against our own interest, because it is material: the
\emph{direction} of the published result is confirmed and its \emph{magnitude} is not, and the
crossover is expected to occur earlier than the 1440p--4K interval measured synthetically. We also
state the caveats bounding the comparison, since real frames differ from the collage in size as well
as in distinctness.

The validation additionally ran on a \emph{different computer}: a four-core Core i7-1165G7 with
15.30~GB of memory, against the eight-core Core i9-11900H with 30.99~GB used for the controlled benchmark, with
the MongoDB cache scaled proportionally from 30~GB to 14~GB. This is therefore a second-hardware
replication and not merely a second-payload one. Absolute figures differ between the machines, as
they must; the architectural ordering does not. A new table in the manuscript specifies both
environments side by side. We also state explicitly that the deployment supplied the cameras and the
capture regime, while the storage measurements were executed on these controlled, fully specified
hosts.

\subsection*{3c. Separating structural claims from measured ones}

Two claims are now identified as independent of our hardware entirely: the 16~MB BSON ceiling is a
protocol limit, so no hardware makes bucketed time-series collections able to store 6K frames; and
the ordering of per-byte costs follows from where each architecture places the payload, not from how
fast the host runs.

\subsection*{3d. Partial rebuttal: the single-host baseline is a design choice}

Respectfully, we retain the single-host, single-client design and wish to justify rather than
silently amend it. Multi-client concurrency changes how energy is attributed and introduces
contention noise. It is effectively a different study, not a more realistic version of this one. The
controlled baseline isolates payload handling, the variable under study, from scheduling and
contention confounds, and keeps the comparison fully reproducible. We have strengthened the evidence
for generality by the means above, which we believe answer the Editor's concern more directly than
added concurrency would, since the crossover is driven by payload size, a structural property of
each storage design, rather than by client count. A multi-client evaluation remains identified as
future work.

\section*{4. Limitations in the abstract}

\point{``The abstract should also explicitly acknowledge the main limitations of the research.''}

\reply Accepted and implemented. The abstract now states, in the authors' own words, that the study
is deliberately bounded and names four limitations: the single-host, single-client controlled
baseline rather than a distributed multi-tenant deployment; a limited photographic corpus, with
real-frame validation reaching 1080p; measurement under one national grid intensity, with the
explicit note that the \emph{ranking} is grid-invariant because emissions scale linearly with that
factor; and embodied-carbon and fleet-scale figures being first-order estimates from published
factors rather than site-specific measurements.

\section*{5. Declaration of AI-assisted technologies}

\point{``If artificial intelligence tools were used for language polishing or translation, this
should be transparently declared in accordance with the journal's editorial policies.''}

\reply Declared. A \textbf{Use of AI-assisted technologies} statement has been added to Statements
and Declarations, recording that generative AI tools were used \emph{solely for language
editing}, that is grammar, clarity, and phrasing of author-written text, and that no AI tool was used to
generate, analyse, or interpret the experimental data, produce figures or tables, or draw the
scientific conclusions. The authors reviewed and edited all such output and take full responsibility
for the entire content. AI is not credited as an author.

We took the opportunity to complete the declarations block, which was otherwise thin: Competing
Interests, Data and Code Availability, and Ethics Approval have also been added. The ethics statement
records that the new real-frame corpus involved no human participants and no personal data, was
captured with no other person present, retained no biometric identifier, and that every saved frame
was verified automatically to contain no detectable face before use.

\section*{Summary of changes}

\begin{tabular}{@{}p{0.30\textwidth}p{0.62\textwidth}@{}}
\toprule
Editor's point & Where addressed \\
\midrule
1. Positioning in the literature &
New Discussion subsection ``Positioning against prior work''; Table~1 extended with a column stating what each prior study concludes and what changes here. \\[2pt]
2. Dialogue with this journal &
New Related Work subsection engaging six recent articles from this journal, each used to sharpen a specific methodological choice. \\[2pt]
3. Applicability beyond our setup &
New subsection deriving the mechanism behind the crossover and giving a measurement procedure for locating it elsewhere; new section reporting the real-case deployment, with a controlled comparison that separates payload construction from hardware; claims that do not depend on measurement identified separately. \\[2pt]
4. Limitations in the abstract &
Abstract revised to name four limitations explicitly. \\[2pt]
5. AI declaration &
Statements and Declarations rewritten, adding use of AI-assisted technologies, competing interests, data availability, and ethics approval. \\
\bottomrule
\end{tabular}

\bigskip
\noindent We hope the revised manuscript now meets the standard required to proceed to an Associate
Editor, and we thank the Editor again for guidance that materially improved the work.

\end{document}
